% !TeX encoding = UTF-8

\chapter{结论与展望}

\section{全文总结}

本文针对联邦学习在实际部署中面临的数据统计异构与隐私效用失衡难题，设计了基于梯度一致性度量的自适应聚合框架与分阶段动态隐私预算分配机制，并完成了非凸收敛理论证明与系统实验验证。主要结论如下：
\begin{enumerate}[label=（\arabic*）, leftmargin=2\ccwd, itemsep=2pt]
  \item 引入历史梯度余弦一致性权重能够显著抑制客户端漂移，在不增加额外通信负载的前提下加快收敛速度；
  \item 分阶段动态噪声调度机制将更多隐私预算分配至模型精调阶段，有效平衡了 $(\varepsilon, \delta)$-差分隐私保障与模型泛化精度；
  \item 在多种异构强度与隐私预算配置下，所提方法在测试准确率与通信效率两项指标上均稳定优于对比基线。
\end{enumerate}

\section{研究局限与未来展望}

尽管本文工作取得了预期成效，但在大规模跨设备异步通信以及抵御拜占庭恶意投毒攻击等方面仍有进一步拓展空间。未来工作将探索面向异步联邦环境的自适应隐私会计方法，并结合硬件可信执行环境提升系统整体安全性。
